% SPDX-License-Identifier: AGPL-3.0-or-later
% Commercial license available
% (c) Concepts 1996-2026 Miroslav Sotek. All rights reserved.
% (c) Code 2020-2026 Miroslav Sotek. All rights reserved.
% ORCID: 0009-0009-3560-0851
% Contact: www.anulum.li | protoscience@anulum.li
% scpn-quantum-control -- JOSS paper preview PDF

\documentclass[11pt]{article}

\usepackage[margin=1in]{geometry}
\usepackage{hyperref}
\usepackage{url}
\usepackage{enumitem}

\hypersetup{
  colorlinks=true,
  linkcolor=blue,
  citecolor=blue,
  urlcolor=blue,
}

\title{\texttt{scpn-quantum-control}: reproducible Kuramoto--XY quantum-control workflows with Rust acceleration and hardware artefact packaging}
\author{Miroslav Sotek\\ANULUM, Marbach SG, Switzerland\\ORCID: 0009-0009-3560-0851}
\date{8 May 2026}

\begin{document}
\maketitle

\section*{Summary}

\texttt{scpn-quantum-control} is a research software package for constructing,
simulating, benchmarking, and packaging quantum-control experiments derived from
heterogeneous Kuramoto oscillator networks and their corresponding XY spin
Hamiltonians. The package maps coupling matrices and natural frequencies to
Qiskit Hamiltonians, variational ansatze, Trotter circuits, simulator workflows,
and IBM Quantum hardware execution records. It combines a Python/Qiskit front
end with Rust/PyO3 acceleration for selected hot-path kernels and includes
scripts that regenerate benchmark tables into JSON and CSV artefacts.

The package has been used to produce a companion IBM Heron r2 hardware preprint
and public validation package on parity-sector and excitation-number correlated
decoherence asymmetry in heterogeneous XY circuits. The same repository now
includes generated artefacts for Rust kernel timings, topology-informed ansatz
construction, VQE comparisons, multi-language coupling-matrix benchmarks,
cross-machine CPU runs, and Vertex T4 GPU batched expectation benchmarks.

\section*{Statement of need}

NISQ quantum-simulation studies often contain a gap between the mathematical
model and the artefacts needed to verify a hardware claim. A complete workflow
must specify how the coupling matrix was generated, how the Hamiltonian was
constructed, how circuits were compiled, which classical baselines were used,
which hardware jobs were submitted, and how raw-count dictionaries were
analysed. Without this chain, small changes in ansatz topology, transpilation,
readout processing, or benchmark implementation can dominate the reported
result.

\texttt{scpn-quantum-control} addresses this gap for Kuramoto--XY and SCPN phase
dynamics experiments. Its target users are quantum software researchers, NISQ
experimentalists, and computational physicists who need reproducible pipelines
connecting oscillator-network models to circuits, simulators, hardware runs,
and published artefacts. The software is not presented as a quantum-advantage
engine; rather, it is an auditable workflow for small- and medium-scale hardware
experiments where honest classical baselines and raw data provenance are
essential.

\section*{State of the field}

Qiskit provides the general quantum-circuit, operator, transpilation, and
execution infrastructure used by the package. Variational workflows such as VQE
and hardware-efficient ansatze are also well established. Existing tools,
however, do not provide an end-to-end Kuramoto--XY workflow that joins
oscillator coupling construction, topology-informed ansatz generation,
Rust-accelerated kernels, IBM raw-count packaging, and generated performance
artefacts in one repository.

The contribution of \texttt{scpn-quantum-control} is therefore not to replace
Qiskit, but to specialise it for a research programme in heterogeneous
oscillator networks, XY Hamiltonians, and symmetry-aware NISQ validation. The
package keeps Qiskit as the circuit and operator layer while adding
domain-specific Hamiltonian construction, experiment manifests, integrity
hashing, and benchmark harnesses.

\section*{Quickstart}

The core workflow constructs a coupling matrix, maps it to a Hamiltonian, and
generates a topology-informed ansatz whose entangling graph follows the
non-zero coupling structure:

\begin{verbatim}
from scpn_quantum_control.bridge.knm_hamiltonian import (
    OMEGA_N_16,
    build_knm_paper27,
    knm_to_ansatz,
    knm_to_hamiltonian,
)

n = 4
k_matrix = build_knm_paper27(n)
hamiltonian = knm_to_hamiltonian(k_matrix, OMEGA_N_16[:n])
ansatz = knm_to_ansatz(k_matrix, reps=2)
\end{verbatim}

The same artefacts can then be passed to VQE routines, simulator workflows,
hardware runners, or benchmark scripts. For hardware execution and packaging,
the repository provides runner and manifest utilities that archive job
identifiers, raw counts, circuit metadata, and integrity hashes.

\section*{Software design}

The architecture separates orchestration from hot kernels. Python modules
provide the public research interface: coupling-matrix construction,
Hamiltonian conversion, ansatz generation, VQE workflows, IBM execution
packaging, and analysis scripts. Rust/PyO3 modules provide deterministic
low-overhead kernels for selected performance-sensitive paths. This design
keeps the package usable by Python-based quantum researchers while allowing
kernel-level optimisation where repeated allocation or object dispatch would
otherwise dominate.

The methods-paper benchmark harnesses follow an artefact-first rule: numerical
tables are regenerated from scripts and committed as JSON/CSV summaries rather
than copied manually into prose. The current generated artefacts include:
Rust/core kernel summaries, ansatz construction summaries, VQE aggregate
tables, multi-language coupling-matrix comparisons, local/ML350/Vertex CPU
runs, and a Vertex T4 batched state-vector expectation benchmark. The GPU
benchmark is intentionally separated from the Rust CPU benchmarks because it
targets batched dense linear algebra for classical validation, not scalar
coupling-matrix construction.

The main public workflow modules include \texttt{bridge/knm\_hamiltonian.py}
for Kuramoto--XY matrix and Hamiltonian conversion,
\texttt{control/structured\_ansatz.py} for topology-informed circuit
construction, \texttt{phase/phase\_vqe.py} for VQE-style workflows, and
hardware runner utilities for simulator/IBM execution. The benchmark harnesses
under \texttt{scripts/} regenerate the paper artefacts rather than requiring
manual table editing.

\section*{Research impact statement}

The companion hardware study provides raw-count archives, job identifiers,
SHA256 integrity hashes, reproduction commands, and bounded statistical claims.
The benchmark artefacts generated on 5 May 2026 show that the project is also
usable as a reproducible methods platform: at $n=16$, the self-contained
coupling-matrix benchmark reports Rust/Python median speedups of 44.1x on the
local workstation, 15.5x on the ML350 server, and 18.1x on Vertex AI CPU. A
separate Vertex Tesla T4 run reports CUDA median times of 0.176--0.905 ms for
batched expectation evaluation at $n=8$--$11$. These results support a hybrid
software strategy in which Rust handles low-latency CPU kernels, GPUs handle
batched dense expectation evaluation, and QPU time is reserved for
hardware-noise questions that classical simulation cannot answer directly.

The package fills a narrow but reusable niche: transparent Kuramoto--XY quantum
simulation workflows with benchmarked software paths and hardware artefact
lineage. This makes it useful both as research infrastructure for future SCPN
experiments and as an example of how small-N NISQ studies can publish their
software, raw data, and performance claims without relying on unverifiable
manual tables.

\section*{AI usage disclosure}

AI-assisted tools were used for drafting and editing support. Numerical claims,
artefact provenance, scientific framing, authorship, and final responsibility
were verified and retained by the author.

\section*{References}

\begin{enumerate}[leftmargin=*]
  \item M. Sotek, \emph{\texttt{scpn-quantum-control}: Quantum-Native SCPN Phase Dynamics and Control}, version 0.9.6, Zenodo record 18821929 (2026).
  \item M. Sotek, ``Replicated hardware observation of parity-sector and excitation-number correlated decoherence asymmetry in the XY Hamiltonian,'' preprint manuscript (2026).
  \item M. Sotek, \emph{DLA parity hardware validation package for heterogeneous Kuramoto--XY circuits}, Zenodo record 19098792 (2026).
  \item Qiskit contributors, \emph{Qiskit: An Open-source Framework for Quantum Computing}, Zenodo record 2573505.
  \item A. Peruzzo et al., ``A variational eigenvalue solver on a photonic quantum processor,'' \emph{Nature Communications} 5, 4213 (2014), doi:10.1038/ncomms5213.
  \item A. Kandala et al., ``Hardware-efficient variational quantum eigensolver for small molecules and quantum magnets,'' \emph{Nature} 549, 242--246 (2017), doi:10.1038/nature23879.
  \item Y. Kuramoto, \emph{Chemical Oscillations, Waves, and Turbulence}, Springer (1984), doi:10.1007/978-3-642-69689-3.
\end{enumerate}

\end{document}
